% Part: proof-theory
% Chapter: natural-deduction
% Section: translation-N2i

\documentclass[../../../include/open-logic-section]{subfiles}

\begin{document}

\olfileid{pt}{nat}{tni}

\olsection{\Log{N2} থেকে \Log{G2}-তে অনুবাদ}

\begin{prop}\ollabel{prop:N2-to-G2}
যদি $\Log{N2c}\ (\Log{N2i}) \Proves
\Gamma \Sequent !A$, তাহলে
$\Log{G2c}\ (\Log{G2i}) + \Cut \Proves
\Gamma' \Sequent \Delta$; এখানে
$\Gamma$ থেকে চিহ্নগুলি সরিয়ে পাওয়া
!!{formula}s-এর বহুসেট হলো~$\Gamma'$,
$!A$ যদি~$\lfalse$ না হয় তবে
$\Delta = \{!A\}$, আর তা হলে
$\Delta = \emptyset$।
% BN-SRC-568 TARGET-CORRECTION-BEGIN
\par\smallskip\noindent\textnormal{\small\textbf{উৎস-সংশোধন (BN-SRC-568):} প্রস্তাবনার ফলপদ্ধতিতে উৎসে ভুলে N2i আছে; N2i প্রমাণের অনুবাদের লক্ষ্য G2i, তাই G2c/G2i জোড়া করা হয়েছে।}
% BN-SRC-568 TARGET-CORRECTION-END
\end{prop}

\begin{proof}
আবার~$\Gamma \Sequent !A$-এর
!!{proof}~$\delta$-র উচ্চতার উপর আরোহ করি।
$\delta$ যদি $x:!A \Sequent !A$ স্বতঃসিদ্ধ
হয়, তাহলে $\pi$ সংশ্লিষ্ট
$!A \Sequent !A$ স্বতঃসিদ্ধ; আর
$!A \ident \lfalse$ হলে
$\lfalse \Sequent \ $ স্বতঃসিদ্ধ।

$\delta$-র শেষে অনুমানবিধি থাকলে
ক্ষেত্র আলাদা করি:
\begin{enumerate}
  \item $R$ যদি \Intro{\lif} হয়,
  প্রধান সূত্র $!B \lif !C$ এবং $x:!B$
  পূর্বধারণার প্রসঙ্গে থাকলেও সিদ্ধান্তের
  প্রসঙ্গে না থাকে, তাহলে $\delta$-র শেষাংশ
  \[
  \AxiomC{}
  \RightLabel{$\delta_1$}
  \Deduce$x:!B, \Gamma \fCenter !C$
  \RightLabel{\Intro{\lif}}
  \UnaryInf$\Gamma \fCenter !B \lif !C$
  \DisplayProof
  \]
  আরোহের অনুমানে $!B, \Gamma'
  \Sequent !C$-এর !!a{proof}~$\pi_1$
  পাই; $!C \ident \lfalse$ হলে এর বদলে
  $!B, \Gamma' \Sequent \ $-এর প্রমাণ পাই।
  পরের ক্ষেত্রে আগে \RightR{\Weakening}
  প্রয়োগ করে~$!C$ ডান দিকে আনি, তারপর
  \RightR{\lif} প্রয়োগে $\pi$ পাই।
  \item $R$ যদি \Intro{\lif} হয় কিন্তু
  $x:!B$ পূর্বধারণার প্রসঙ্গে না থাকে,
  তাহলে $\delta$-র শেষাংশ
  \[
  \AxiomC{}
  \RightLabel{$\delta_1$}
  \Deduce$\Gamma \fCenter !C$
  \RightLabel{\Intro{\lif}}
  \UnaryInf$\Gamma \fCenter !B \lif !C$
  \DisplayProof
  \]
  আরোহের অনুমানে $\Gamma' \Sequent !C$-এর
  !!a{proof}~$\pi_1$ পাই। যদি
  $!C \ident \lfalse$ হয়, তবে আরোহে
  ডান অংশ ফাঁকা আসে; আগে
  \RightR{\Weakening} দিয়ে~$!C$ ফিরিয়ে
  নিই। তারপর $!B$-সহ \LeftR{\Weakening}
  এবং শেষে \RightR{\lif} প্রয়োগ করে
  $\pi$ পাই।
  % BN-SRC-569 TARGET-CORRECTION-BEGIN
  \par\smallskip\noindent\textnormal{\small\textbf{উৎস-সংশোধন (BN-SRC-569):} অবমুক্তিযোগ্য $x:!B$ অনুপস্থিত ক্ষেত্রেও $!C=\lfalse$ হতে পারে; উৎসে ফাঁকা ডান অংশ থেকে →-ভূমিকার আগে \RightR{\Weakening} ধাপটি বাদ।}
  % BN-SRC-569 TARGET-CORRECTION-END
  \item $R$ যদি \Elim{\lif} হয়,
  $\delta$-র শেষাংশ
  \[
  \AxiomC{}
  \RightLabel{$\delta_1$}
  \Deduce$\Gamma_1 \fCenter !B \lif !A$
  \AxiomC{}
  \RightLabel{$\delta_2$}
  \Deduce$\Gamma_2 \fCenter !B$
  \RightLabel{\Elim{\lif}}
  \BinaryInf$\Gamma_1, \Gamma_2 \fCenter !A$
  \DisplayProof
  \]
  যেখানে $\Gamma = \Gamma_1 \cup \Gamma_2$।
  আরোহের অনুমানে $\Gamma_1' \Sequent
  !B \lif !A$-এর !!{proof}s~$\pi_1$
  এবং $\Gamma_2' \Sequent !B$-এর
  $\pi_2$ পাই। $!B \ident \lfalse$ হলে
  দ্বিতীয় আরোহফলের ডান অংশ ফাঁকা;
  \RightR{\Weakening} দিয়ে~$!B$ এনে
  তবেই নিচের গঠন নিই। !!{proof}~$\pi$:
  \[
  \AxiomC{}
  \RightLabel{$\pi_1$}
  \Deduce$\Gamma_1' \fCenter !B \lif !A$
  \AxiomC{}
  \RightLabel{$\pi_2$}
  \Deduce$\Gamma_2' \fCenter !B$
  \Axiom$!A \fCenter !A$
  \RightLabel{\LeftR{\lif}}
  \BinaryInf$!B \lif !A, \Gamma_2' \fCenter !A$
  \RightLabel{\Cut}
  \BinaryInf$\Gamma_1', \Gamma_2' \fCenter !A$
  \DisplayProof
  \]
  বহুসেট $\Gamma_1' \cup \Gamma_2'$
  সব সময় $\Gamma' =
  (\Gamma_1 \cup \Gamma_2)'$-এর সমান নয়:
  যেমন, $\Gamma_1$ ও $\Gamma_2$ উভয়ে
  $x:!C$ থাকলে $\Gamma_1' \cup \Gamma_2'$-তে $!C$-এর
  দুটি অনুলিপি থাকে, কিন্তু
  $(\Gamma_1 \cup \Gamma_2)'$-তে থাকে
  মাত্র একটি। উপযুক্ত \LeftR{\Contraction}
  প্রয়োগে এটি সংশোধন করা যায়।
  % BN-SRC-570 TARGET-CORRECTION-BEGIN
  \par\smallskip\noindent\textnormal{\small\textbf{উৎস-সংশোধন (BN-SRC-570):} দ্বিতীয় উপপ্রমাণের অন্তিম সূত্র $!B$ নিজেই $\lfalse$ হতে পারে; তখন আরোহফলের ফাঁকা ডান অংশে আগে ডান-দুর্বলীকরণ দরকার। উৎসের সাধারণ বাক্যটি এই ক্ষেত্র বাদ দিয়েছে।}
  % BN-SRC-570 TARGET-CORRECTION-END

  $!A \ident \lfalse$ হলে উপরের
  !!{proof}-এ প্রথমে
  $\Gamma_1', \Gamma_2' \Sequent \lfalse$
  পাই। লক্ষ্য ডান অংশ ফাঁকা; তাই
  $\lfalse \Sequent \ $ স্বতঃসিদ্ধের
  সঙ্গে~$\lfalse$-এ \Cut প্রয়োগ করে
  $\Gamma_1', \Gamma_2' \Sequent \ $
  পাই, তারপর দরকারে বাঁ-সংকোচন করে
  লক্ষ্য !!a{proof}~$\pi$ নিই।
  % BN-SRC-571 TARGET-CORRECTION-BEGIN
  \par\smallskip\noindent\textnormal{\small\textbf{উৎস-সংশোধন (BN-SRC-571):} উৎসে $!A=\lfalse$ হলে ভুল করে $\delta_2$ থেকে ফাঁকা ডান অংশ দাবি করা হয়েছে; $\delta_2$-র অন্তিম সূত্র $!B$, যা মিথ্যা হওয়া বাধ্যতামূলক নয়। উপরের →-নিষ্কাশন বৃক্ষ থেকে $\lfalse$ পেয়ে মিথ্যা-স্বতঃসিদ্ধের সঙ্গে কর্তনে ফাঁকা ডান অংশ গঠিত।}
  % BN-SRC-571 TARGET-CORRECTION-END
\item $R$ যদি \FalseCl হয়,
  $\delta$-র শেষাংশ
  \[
  \AxiomC{}
  \RightLabel{$\delta_1$}
  \Deduce$x: \lnot !A, \Gamma \fCenter \lfalse$
  \RightLabel{\FalseCl}
  \UnaryInf$\Gamma \fCenter !A$
  \DisplayProof
  \]
  আরোহের অনুমানে $\lnot !A,
  \Gamma' \Sequent \ $-এর
  !!a{proof}~$\pi_1$ আছে। !!{proof}~$\pi$
  নিই
  \[
  \Axiom$!A \fCenter !A$
  \RightLabel{\RightR{\lnot}}
  \UnaryInf$\fCenter !A, \lnot !A$
  \AxiomC{}
  \RightLabel{$\pi_1$}
  \Deduce$\lnot !A, \Gamma' \fCenter $
  \RightLabel{\Cut}
  \BinaryInf$\Gamma' \fCenter !A$
  \DisplayProof
  \]
  $!A \ident \lfalse$ হলে এখানেও
  $\lfalse \Sequent \ $ স্বতঃসিদ্ধের
  সঙ্গে শেষ~$\lfalse$-এ \Cut প্রয়োগে
  প্রয়োজনীয় ফাঁকা ডান অংশ পাই।
  % BN-SRC-572 TARGET-CORRECTION-BEGIN
  \par\smallskip\noindent\textnormal{\small\textbf{উৎস-সংশোধন (BN-SRC-572):} G2 প্রমাণবৃক্ষের দ্বিতীয় পূর্বধারণায় উৎসের N2 চিহ্ন $x:$ সরানো হয়েছে; G2 বহুসেটে চিহ্ন নেই। $!A=\lfalse$ হলে প্রস্তাবনার ফাঁকা ডান অংশ পেতে মিথ্যা-স্বতঃসিদ্ধের সঙ্গে একটি অতিরিক্ত কর্তনও স্পষ্ট করা হয়েছে।}
  % BN-SRC-572 TARGET-CORRECTION-END
\end{enumerate}
শুধু \FalseCl{}-এর অনুবাদেই
!!a{proof}~$\pi$-র কোনো সিকোয়েন্টের
ডান দিকে একাধিক !!{formula} আসে।
অর্থাৎ \Log{N2i}-এর !!{proof}-এর
অনুবাদ \Log{G2i}+\Cut-!!{proof}।
% BN-SRC-573 TARGET-CORRECTION-BEGIN
\par\smallskip\noindent\textnormal{\small\textbf{উৎস-সংশোধন (BN-SRC-573):} উৎসের শেষ বাক্যে কর্তনহীন G2i-প্রমাণ বলা হলেও →-নিষ্কাশনের প্রদর্শিত অনুবাদেই \Cut আছে; প্রস্তাবনা অনুযায়ী G2i+\Cut বলা হয়েছে।}
% BN-SRC-573 TARGET-CORRECTION-END
\end{proof}

\end{document}
